
\documentclass[10pt,a4paper]{article}
\usepackage[margin=0.65in]{geometry}
\usepackage[T1]{fontenc}
\usepackage[utf8]{inputenc}
\usepackage{lmodern}
\usepackage{microtype}
\usepackage{multicol}
\usepackage{booktabs}
\usepackage{array}
\usepackage[table]{xcolor}
\usepackage{graphicx}
\usepackage{tikz}
\usepackage{hyperref}
\setlength{\parindent}{0pt}
\setlength{\parskip}{0.35em}
\pagestyle{empty}
\definecolor{accent}{named}{black}
\begin{document}
\begin{center}
{\LARGE\bfseries Applied Optics}\par
\vspace{0.25em}
{\small Journal Manuscript Preview Library}\par
\vspace{0.15em}
{\footnotesize Optics-oriented two-column journal format}\par
\vspace{0.2em}
{\footnotesize Profile slug: applied-optics \quad Family: optica-publishing \quad Scope: Specific journal profile}
\end{center}
\vspace{0.2em}



\noindent\textbf{Authors.} A. Placeholder, B. Example, and C. Template\par
\noindent\textbf{Affiliations.} Department of Simulated Formatting, Placeholder Institute\par
\noindent\textbf{Contact.} preview@example.org\par
\vspace{0.4em}
\noindent\textbf{Abstract.} This simulated first-page preview uses placeholder text, a placeholder figure, and a placeholder table to visualize how a manuscript may look under the selected journal profile. The content is intentionally generic so attention stays on spacing, title treatment, abstract presentation, caption location, and the overall page rhythm.\par\vspace{0.6em}
\noindent\textbf{Keywords:} placeholder manuscript, journal layout, figure slot, table slot, references.\par\vspace{0.6em}

\begin{multicols}{2}
\small
\textbf{1. Introduction.} This placeholder article body is intentionally short but dense enough to reveal the page rhythm of the selected journal profile. It emphasizes title treatment, abstract position, section spacing, figure placement, table styling, and bibliography preview while keeping the manuscript content secondary.

\textbf{2. Methods.} The simulated manuscript uses generic technical prose, a compact table, and an inline figure block so that the resulting PNG can be reviewed visually as a format check rather than as a scientific contribution preview.


\begin{center}\fcolorbox{black!60}{gray!10}{\begin{minipage}[c][2.25cm][c]{0.93\linewidth}\centering\Large Image Placeholder\\[0.4em]\normalsize Simulated figure area\end{minipage}}\end{center}\noindent\textbf{Figure 1.} Placeholder figure caption showing the expected caption style and relative spacing.\par\vspace{0.5em}

\textbf{3. Results.} Placeholder narrative text continues here to show section headings, paragraph spacing, and how a figure or table sits relative to adjacent content in the current preview style.

\noindent\textbf{Table 1.} Placeholder table caption showing metric labels and compact journal spacing.\par\begin{center}
\begin{tabular}{lccc}
\toprule
Method & Score & Error & Time\\
\midrule
Placeholder A & 0.91 & 0.08 & 1.4\\
Placeholder B & 0.88 & 0.11 & 1.7\\
Placeholder C & 0.84 & 0.15 & 2.0\\
\bottomrule
\end{tabular}\end{center}\vspace{0.5em}

\textbf{4. Discussion.} The content is arbitrary by design. The main goal is to provide a quick visual reference for manuscript structure, front matter, floats, and house-style signals inside the profile library.

\end{multicols}

\vspace{0.5em}
\noindent \textbf{References.} [1] A. Smith and B. Lee, ``Journal template adaptation,'' \textit{Placeholder Journal}, 2024. [2] C. Garcia \textit{et al.}, ``Figure and table policy notes,'' 2025.\par
\vspace{0.6em}
\fcolorbox{black!45}{gray!6}{\begin{minipage}{0.98\linewidth}
\footnotesize Simulated preview built from the local starter profile library.\\ Use the official template to finalize submission-ready formatting.
\end{minipage}}
\end{document}
